%%%%%%%%%%%%%%%%%%%%%%%%%%%%%%%%%%%%%%%%%%%%%%%%%%%%%%%%%%%%%%%%%%%%%%
\section{The Albert Algebra and Exceptional Structure}
\label{sec:albert:main}
%%%%%%%%%%%%%%%%%%%%%%%%%%%%%%%%%%%%%%%%%%%%%%%%%%%%%%%%%%%%%%%%%%%%%%

This section assembles the algebraic foundation on which the whole
construction rests: the exceptional Jordan (Albert) algebra
$J_{3}(\mathbb{O})$, its trace and cubic invariants, the sharp map, the
notion of rank, and the four real forms of the exceptional groups that
act on it. We isolate, with care, \emph{which} invariant is preserved by
\emph{which} group, since the entire vacuum-selection argument
(\autoref{sec:potential:main}) and the Lorentzian soldering
(\autoref{sec:soldering:main}) depend on getting this exactly right. The
heavy octonionic identities (associator gymnastics, the explicit cubic
norm in components, the Freudenthal cross product) are deferred to
\autoref{app:albert:main}; here we state the structural facts and prove
the ones that are short. Throughout, the metric signature is
$(-,+,+,+)$, $[f]=\mathrm{mass}$, and the dimensionless couplings are
$\hat g^{2},\xi,\hat y^{2},\kappa$.

%%%%%%%%%%%%%%%%%%%%%%%%%%%%%%%%%%%%%%%%%%%%%%%%%%%%%%%%%%%%%%%%%%%%%%
\subsection{The exceptional Jordan algebra \texorpdfstring{$J_{3}(\mathbb{O})$}{J3(O)}}
\label{sec:albert:def}
%%%%%%%%%%%%%%%%%%%%%%%%%%%%%%%%%%%%%%%%%%%%%%%%%%%%%%%%%%%%%%%%%%%%%%

Let $\mathbb{O}$ denote the (real, division) octonions, with conjugation
$x\mapsto \bar x$, norm $n(x)=x\bar x\in\mathbb{R}_{\ge 0}$, and real part
$\operatorname{Re} x=\tfrac12(x+\bar x)$. Octonionic multiplication is
nonassociative; this nonassociativity is exactly what forces the
\emph{exceptional} structure below and obstructs any matrix-algebra
shortcut.

\begin{definition}[Albert algebra]\label{def:albert:J3O}
The \emph{Albert algebra} is the space of $3\times 3$ octonionic
Hermitian matrices
\begin{equation}\label{eq:albert:J3Odef}
  J_{3}(\mathbb{O})
  =\Bigl\{\,X=X^{\dagger}\in M_{3}(\mathbb{O})\,\Bigr\}
  =\left\{
  \begin{pmatrix}
   \xi_{1} & z_{3} & \bar z_{2}\\
   \bar z_{3} & \xi_{2} & z_{1}\\
   z_{2} & \bar z_{1} & \xi_{3}
  \end{pmatrix}
  : \xi_{i}\in\mathbb{R},\; z_{i}\in\mathbb{O}\right\},
\end{equation}
equipped with the commutative \emph{Jordan product}
\begin{equation}\label{eq:albert:jordanprod}
  X\circ Y=\tfrac12\,(XY+YX).
\end{equation}
\end{definition}

The three real diagonal entries and three octonionic off-diagonal
entries give
\begin{equation}\label{eq:albert:dim}
  \dim_{\mathbb{R}} J_{3}(\mathbb{O})
  = 3\cdot 1 + 3\cdot 8 = 27 .
\end{equation}

\begin{theorem}[Formally real Jordan algebra]\label{thm:albert:formallyreal}
$\bigl(J_{3}(\mathbb{O}),\circ\bigr)$ is a $27$-dimensional, commutative,
formally real (Euclidean) Jordan algebra: the Jordan identity
$X\circ\bigl(Y\circ X^{\circ 2}\bigr)=\bigl(X\circ Y\bigr)\circ X^{\circ 2}$
holds, and $\sum_{i} X_{i}\circ X_{i}=0$ implies every $X_{i}=0$. It is
the unique exceptional simple formally real Jordan algebra (it is not
special, i.e.\ not a Jordan subalgebra of any associative algebra under
the symmetrised product).
\end{theorem}

\begin{proof}[Proof sketch]
Commutativity is manifest from \eqref{eq:albert:jordanprod}. The Jordan
identity is a power-associativity statement that, by linearisation,
reduces to relations involving at most two octonions at a time; since any
two octonions generate an associative (indeed commutative) subalgebra
$\cong\mathbb{C}$ or $\mathbb{H}$ (Artin's theorem), the identity holds
despite global nonassociativity. Formal reality follows from positivity
of the trace form \eqref{eq:albert:traceform} below. Exceptionality
(non-speciality) is the Albert--Jacobson theorem; see
\autoref{app:albert:main} and \cite{Albert:1934,Jacobson:1968,McCrimmon:2004}.
\end{proof}

%%%%%%%%%%%%%%%%%%%%%%%%%%%%%%%%%%%%%%%%%%%%%%%%%%%%%%%%%%%%%%%%%%%%%%
\subsection{Trace form, cubic norm, and the sharp map}
\label{sec:albert:invariants}
%%%%%%%%%%%%%%%%%%%%%%%%%%%%%%%%%%%%%%%%%%%%%%%%%%%%%%%%%%%%%%%%%%%%%%

Three invariant objects organise the algebra: a quadratic trace form, a
cubic norm, and a quadratic-in-$X$ \emph{sharp} (adjoint) map. Their
distinct invariance groups are the crux of the whole construction.

\paragraph{Linear and quadratic traces.}
Define the linear trace and the symmetric \emph{trace form}
\begin{align}
  \operatorname{tr}(X) &= \xi_{1}+\xi_{2}+\xi_{3},
  \label{eq:albert:lintrace}\\[2pt]
  \langle X,Y\rangle &= \operatorname{Tr}(X\circ Y),
  \qquad
  Q(X)\equiv \operatorname{Tr}(X^{2})=\langle X,X\rangle .
  \label{eq:albert:traceform}
\end{align}
In components,
\begin{equation}\label{eq:albert:Qcomp}
  Q(X)=\operatorname{Tr}(X^{2})
  =\sum_{i}\xi_{i}^{2}+2\sum_{i} n(z_{i}) \;\ge 0 ,
\end{equation}
so $Q$ is positive-definite: $J_{3}(\mathbb{O})$ is Euclidean, and $Q$ is
the natural $\mathrm{O}(27)$-type inner product. This positivity is what
makes the dynamical coset metric definite later
(\autoref{sec:coset:main}).

\paragraph{Cubic norm.}
The \emph{cubic norm} (the octonionic ``determinant'') is
\begin{equation}\label{eq:albert:Ndef}
  N(X)=\det\nolimits_{\mathbb{O}}(X)
  =\xi_{1}\xi_{2}\xi_{3}
   -\sum_{i}\xi_{i}\,n(z_{i})
   +2\,\operatorname{Re}\!\bigl(z_{1}z_{2}z_{3}\bigr).
\end{equation}
The bracketing of the cubic term $z_{1}z_{2}z_{3}$ is unambiguous under
$\operatorname{Re}(\cdot)$ despite nonassociativity (a standard
octonionic identity; \autoref{app:albert:main}). Unlike $Q$, the cubic
norm is \emph{odd} under $X\mapsto -X$ and is genuinely indefinite; its
restriction to the $H_{2}(\mathbb{O})$ sub-block supplies the Minkowski
form used in soldering (\autoref{sec:soldering:main}).

\paragraph{Sharp map.}
The \emph{sharp} (adjoint) map $X\mapsto X^{\#}$ is the gradient of $N$
with respect to the trace form,
\begin{equation}\label{eq:albert:sharpdef}
  \langle X^{\#}, Y\rangle = \operatorname{D}N(X)[Y]
  \quad\Longleftrightarrow\quad
  X^{\#}=\nabla N(X),
\end{equation}
equivalently the quadratic map characterised by
\begin{equation}\label{eq:albert:sharpid}
  X\circ X^{\#}=N(X)\,\mathbf 1,
  \qquad
  (X^{\#})^{\#}=N(X)\,X .
\end{equation}
In a \emph{Jordan frame} $X=\operatorname{diag}(\lambda_{1},\lambda_{2},\lambda_{3})$,
\begin{equation}\label{eq:albert:sharpframe}
  X^{\#}=\operatorname{diag}\!\bigl(\lambda_{2}\lambda_{3},\,
  \lambda_{3}\lambda_{1},\,\lambda_{1}\lambda_{2}\bigr),
  \qquad
  \langle X^{\#},X^{\#}\rangle
  =(\lambda_{2}\lambda_{3})^{2}+(\lambda_{3}\lambda_{1})^{2}+(\lambda_{1}\lambda_{2})^{2}\ge 0 .
\end{equation}
The combination $\langle X^{\#},X^{\#}\rangle$ is a sum of squares,
hence non-negative; this is the quantity whose vanishing characterises
the vacuum orbit.

%%%%%%%%%%%%%%%%%%%%%%%%%%%%%%%%%%%%%%%%%%%%%%%%%%%%%%%%%%%%%%%%%%%%%%
\subsection{Rank, idempotents, and the rank-one criterion}
\label{sec:albert:rank}
%%%%%%%%%%%%%%%%%%%%%%%%%%%%%%%%%%%%%%%%%%%%%%%%%%%%%%%%%%%%%%%%%%%%%%

\begin{definition}[Rank]\label{def:albert:rank}
The \emph{rank} of $X\in J_{3}(\mathbb{O})$ is the degree of its minimal
polynomial; equivalently, in a Jordan frame it is the number of nonzero
eigenvalues $\lambda_{i}$. Thus $\operatorname{rank}(X)\in\{0,1,2,3\}$,
with $\operatorname{rank}(0)=0$. A \emph{primitive idempotent} is a
rank-one $E$ with $E\circ E=E$ (equivalently $E^{2}=E$, $\operatorname{tr}E=1$).
\end{definition}

\begin{theorem}[Rank-one criterion]\label{thm:albert:rankone}
For all $X\in J_{3}(\mathbb{O})$,
\begin{equation}\label{eq:albert:rankone}
  \operatorname{rank}(X)\le 1
  \quad\Longleftrightarrow\quad
  X^{\#}=0 .
\end{equation}
Moreover $\operatorname{rank}(X)\le 2\Leftrightarrow N(X)=0$ (with
$X^{\#}\neq 0$), and $\operatorname{rank}(X)=3\Leftrightarrow N(X)\neq 0$.
\end{theorem}

\begin{proof}
Every $X$ is conjugate by $F_{4}=\operatorname{Aut}(J_{3}(\mathbb{O}))$
(\autoref{thm:albert:fourforms}) to a Jordan frame
$\operatorname{diag}(\lambda_{1},\lambda_{2},\lambda_{3})$
(spectral theorem for Euclidean Jordan algebras). By
\eqref{eq:albert:sharpframe}, $X^{\#}=0$ iff all three pairwise products
$\lambda_{i}\lambda_{j}$ vanish, i.e.\ at most one $\lambda_{i}$ is
nonzero, i.e.\ $\operatorname{rank}(X)\le1$. The norm statements follow
from $N(X)=\lambda_{1}\lambda_{2}\lambda_{3}$ in a frame: $N=0$ iff some
$\lambda_{i}=0$ (rank $\le 2$), and $N\neq0$ iff all nonzero (rank $3$).
Since $F_{4}$ preserves both $N$ and $X^{\#}$, the criteria are
frame-independent.
\end{proof}

\begin{remark}[Scope]\label{rmk:albert:scope}
\autoref{thm:albert:rankone} is a \emph{theorem} of Jordan algebra, used
verbatim in the vacuum-selection argument: the $E_{6}$-covariant term
$\langle X^{\#},X^{\#}\rangle$ in the potential $V$ vanishes
\emph{exactly} on the rank-$\le1$ cone (\autoref{sec:potential:main}).
The selection of the rank-one orbit as the global minimum is proved
there; here we only supply the algebraic criterion.
\end{remark}

%%%%%%%%%%%%%%%%%%%%%%%%%%%%%%%%%%%%%%%%%%%%%%%%%%%%%%%%%%%%%%%%%%%%%%
\subsection{The four real forms and what each preserves}
\label{sec:albert:groups}
%%%%%%%%%%%%%%%%%%%%%%%%%%%%%%%%%%%%%%%%%%%%%%%%%%%%%%%%%%%%%%%%%%%%%%

The single most consequential fact of this section is that
\emph{different groups preserve different invariants}. Conflating them
(in particular, asserting that the structure group preserves the trace
form $Q$, or that $(Q,N)$ separates orbits) is the root error against
which the entire construction is hardened.

\begin{theorem}[Four real forms]\label{thm:albert:fourforms}
The following groups act on $J_{3}(\mathbb{O})$ and the real $27$:
\begin{enumerate}
\item\label{it:albert:F4}
  $F_{4}=\operatorname{Aut}\bigl(J_{3}(\mathbb{O}),\circ\bigr)$ is the
  \textbf{compact} exceptional group, $\dim F_{4}=52$. It preserves the
  full Jordan product, hence \emph{both} the trace form
  $Q=\operatorname{Tr}(X^{2})$ \emph{and} the cubic norm $N$, and acts as
  $F_{4}\subset\mathrm{SO}(27)$ (it is the stabiliser of the identity
  $\mathbf 1$).
\item\label{it:albert:E6m26}
  The \textbf{reduced structure group}
  $E_{6(-26)}$ is the \textbf{non-compact} real form,
  $\dim E_{6(-26)}=78$. It preserves the cubic norm $N$ \emph{up to a
  real positive character} $\lambda(g)$,
  \begin{equation}\label{eq:albert:Ncharacter}
    N(g\cdot X)=\lambda(g)\,N(X),\qquad \lambda:E_{6(-26)}\to\mathbb{R}_{>0},
  \end{equation}
  and does \emph{not} preserve $Q$. (If it did, it would lie in
  $\mathrm{SO}(27)$ and be compact, a contradiction.)
\item\label{it:albert:E6c}
  The \textbf{compact} real form, here denoted compact $E_{6}$,
  $\dim E_{6}=78$, acts unitarily on the complexified $27_{\mathbb{C}}$;
  on it the norm character is identically $1$ (a continuous homomorphism
  from a compact connected group to $\mathbb{R}_{>0}$ is trivial), and
  $Q$ is the invariant Hermitian form.
\item\label{it:albert:E6m14}
  The real form $E_{6(-14)}$, $\dim=78$, has \textbf{maximal compact
  subgroup} $\operatorname{Spin}(10)\times\mathrm{U}(1)$.
\end{enumerate}
\end{theorem}

\begin{proof}[Proof sketch]
\ref{it:albert:F4} An automorphism of the Jordan product preserves
powers, hence the trace form (via $Q(X)=\operatorname{Tr}X^{2}$) and the
norm (via Cayley--Hamilton in the algebra); compactness and $\dim 52$ are
classical \cite{Springer:2000}. \ref{it:albert:E6m26} The reduced
structure group is by definition the norm-similarity group;
\eqref{eq:albert:Ncharacter} is its defining property, and non-compactness
follows because a non-trivial similarity character exists. \ref{it:albert:E6c}
Triviality of the character on a compact group is the homomorphism
argument given. \ref{it:albert:E6m14} is the standard real-form
classification; the maximal compact of $E_{6(-14)}$ is
$\operatorname{Spin}(10)\times\mathrm{U}(1)$
\cite{Springer:2000,Slansky:1981,Helgason:1978}.
\end{proof}

\begin{postulate}[Real form for dynamics]\label{post:albert:realform}
For the \emph{dynamics} of UCFT we commit to the \textbf{compact} real
form (equivalently $E_{6(-14)}$, whose maximal compact is exactly
$H=\operatorname{Spin}(10)\times\mathrm{U}(1)$). Then $H$ is compact, the
coset metric on $E_{6}/H$ is positive-definite, the kinetic term is
healthy, and the gauge sector is unitary. The non-compact form
$E_{6(-26)}$ is used \textbf{only} as the rigid orbit-classifier of the
real $27$. This choice is part of the postulate set (it is the form on
which the potential $V$ is coercive; non-compact orbits cannot be
coercive, since the norm character \eqref{eq:albert:Ncharacter} is then
unbounded). Cf.\ \hyperref[post:potential:P1]{P1} and the vacuum theorem
in \autoref{sec:potential:main}.
\end{postulate}

\begin{theorem}[Orbit classification by rank and norm]\label{thm:albert:orbits}
The orbits of $E_{6(-26)}$ on the real $27=J_{3}(\mathbb{O})$ are
classified by the pair $\bigl(\operatorname{rank}(X),\,N(X)\bigr)$, and
\emph{not} by $\bigl(Q(X),N(X)\bigr)$. Explicitly, the nonzero orbit
strata are
\begin{equation}\label{eq:albert:orbitstrata}
  \underbrace{\operatorname{rank}=1}_{X^{\#}=0}\;\subset\;
  \underbrace{\operatorname{rank}=2}_{X^{\#}\neq 0,\,N=0}\;\subset\;
  \underbrace{\operatorname{rank}=3}_{N\neq 0}\,,
\end{equation}
the rank-$3$ stratum further refined by the sign / value of $N$. The
trace form $Q$ is \textbf{not} an $E_{6}$-invariant (only an
$F_{4}$-invariant), so any ``$(Q,N)$ separates orbits'' statement is
false.
\end{theorem}

\begin{proof}
$E_{6(-26)}$ preserves $N$ up to the character
\eqref{eq:albert:Ncharacter} and acts transitively on each rank stratum
of fixed $N$-class (Springer; \autoref{app:albert:main}). The rank is an
algebraic invariant (\autoref{def:albert:rank}), preserved because the
sharp map is $E_{6}$-covariant: $(g\cdot X)^{\#}=\lambda(g)\,
(g^{\#})\!\cdot X^{\#}$ in the appropriate contragredient action, so
$X^{\#}=0$ is a group-invariant condition (and likewise $N=0$). That $Q$
is not invariant is \autoref{thm:albert:fourforms}\ref{it:albert:E6m26}:
a structure-group element with $\lambda(g)\neq1$ generically rescales
$\operatorname{Tr}(X^{2})$.
\end{proof}

\begin{corollary}[Closed minimal orbit]\label{cor:albert:minorbit}
The unique closed nonzero $E_{6}$-orbit in $\mathbb{P}(27_{\mathbb{C}})$
is the rank-one (primitive-idempotent) orbit, the \emph{complex Cayley
plane}
\begin{equation}\label{eq:albert:EIII}
  \mathcal{M}
  = E_{6}/\bigl(\operatorname{Spin}(10)\times\mathrm{U}(1)\bigr)
  = \mathrm{EIII},
  \qquad
  \dim_{\mathbb{R}}\mathcal{M}=32 .
\end{equation}
This is the orbit of a rank-one primitive idempotent, characterised by
$X^{\#}=0$ (\autoref{thm:albert:rankone}); it is \textbf{not} the orbit
of $\operatorname{diag}(1,1,1)$ (which is rank $3$, has
$\langle X^{\#},X^{\#}\rangle=3>0$, stabiliser $F_{4}$, orbit $E_{6}/F_{4}$
of real dimension $26$).
\end{corollary}

\begin{proof}
The closed orbit of a reductive group in a projective space is the orbit
of a highest-weight line; for the $27$ of $E_{6}$ this is the rank-one
line, with stabiliser the parabolic whose reductive part fixes the
$H$-decomposition \eqref{eq:albert:adjbranch}. Over the compact form the
isotropy is $H=\operatorname{Spin}(10)\times\mathrm{U}(1)$, and
$\dim_{\mathbb{R}}E_{6}-\dim_{\mathbb{R}}H=78-46=32$. The real rank-one
idempotents alone form $F_{4}/\operatorname{Spin}(9)\cong
\mathbb{OP}^{2}$ of dimension $16$; the full complex line under $E_{6}$
sweeps the larger $32$-dimensional $\mathrm{EIII}$
\cite{Helgason:1978,Springer:2000}.
\end{proof}

\begin{remark}[Consistency constraints]\label{rmk:albert:forbidden}
Three statements are \emph{false} and are never used:
(i) that the structure group preserves $Q=\operatorname{Tr}(X^{2})$
(only $F_{4}$ does);
(ii) that $(Q,N)$ separates $E_{6}$-orbits (orbits are classified by rank
and $N$, \autoref{thm:albert:orbits});
(iii) that the $32$-dimensional vacuum is the orbit of
$\operatorname{diag}(1,1,1)$ with stabiliser $F_{4}$ (that orbit is
$E_{6}/F_{4}$, dimension $26$; the vacuum is the rank-one EIII orbit,
\autoref{cor:albert:minorbit}).
\end{remark}

%%%%%%%%%%%%%%%%%%%%%%%%%%%%%%%%%%%%%%%%%%%%%%%%%%%%%%%%%%%%%%%%%%%%%%
\subsection{Branchings under \texorpdfstring{$H=\operatorname{Spin}(10)\times\mathrm{U}(1)$}{H = Spin(10) x U(1)}}
\label{sec:albert:branchings}
%%%%%%%%%%%%%%%%%%%%%%%%%%%%%%%%%%%%%%%%%%%%%%%%%%%%%%%%%%%%%%%%%%%%%%

The isotropy group $H=\operatorname{Spin}(10)\times\mathrm{U}(1)$ of the
vacuum orbit \eqref{eq:albert:EIII} is exactly the gauge group of the
construction. The relevant representations branch as follows; these are
the canonical decompositions used everywhere downstream.

\begin{theorem}[Defining branchings]\label{thm:albert:branchings}
Under $H=\operatorname{Spin}(10)\times\mathrm{U}(1)$ (subscripts are the
$\mathrm{U}(1)$ charges),
\begin{align}
  \mathbf{78}
   &= \mathbf{45}_{0}\oplus\mathbf{1}_{0}
      \oplus\mathbf{16}_{-3}\oplus\overline{\mathbf{16}}_{+3}
      \qquad\text{(adjoint),}
   \label{eq:albert:adjbranch}\\[2pt]
  \mathbf{27}
   &= \mathbf{1}_{+4}\oplus\mathbf{10}_{-2}\oplus\mathbf{16}_{+1}
      \qquad\text{(fundamental / matter).}
   \label{eq:albert:fundbranch}
\end{align}
Consequently the coset tangent space is
\begin{equation}\label{eq:albert:cosettangent}
  \mathfrak{m}
   = \mathbf{16}_{-3}\oplus\overline{\mathbf{16}}_{+3}
   \qquad(\,32\text{ real dimensions}\,),
\end{equation}
which contains \textbf{no} $\operatorname{Spin}(10)$ singlet.
\end{theorem}

\begin{proof}
The decompositions \eqref{eq:albert:adjbranch}--\eqref{eq:albert:fundbranch}
are the standard $E_{6}\downarrow\operatorname{Spin}(10)\times\mathrm{U}(1)$
branchings \cite{Slansky:1981}; the $\mathrm{U}(1)$ charges are fixed by
requiring the cubic $E_{6}$-invariant $\mathbf{27}^{3}$ and the adjoint
commutators to be charge-neutral. The coset directions are the
non-$H$ generators in \eqref{eq:albert:adjbranch}, namely
$\mathbf{16}_{-3}\oplus\overline{\mathbf{16}}_{+3}$, of real dimension
$16+16=32=\dim_{\mathbb{R}}\mathcal{M}$ (\autoref{cor:albert:minorbit}).
Dimension and charge counting exclude any singlet in $\mathfrak{m}$.
\end{proof}

\begin{theorem}[Irreducibility and the definite coset metric]\label{thm:albert:irrep}
The coset tangent module $\mathfrak{m}=\mathbf{16}_{-3}\oplus
\overline{\mathbf{16}}_{+3}$ is real-irreducible \emph{of complex type}
(its commutant is $\mathbb{C}$). The $H$-invariant metric on
$\mathfrak{m}$ is therefore \textbf{unique up to scale} and, on the
compact real form (\autoref{post:albert:realform}),
\textbf{positive-definite}.
\end{theorem}

\begin{proof}
The $\mathrm{U}(1)$ charge $\pm3$ forbids the symmetric self-pairings
$\mathbf{16}_{-3}\!\otimes\!\mathbf{16}_{-3}$ and
$\overline{\mathbf{16}}_{+3}\!\otimes\!\overline{\mathbf{16}}_{+3}$
(charges $\mp6\neq0$), leaving only the charge-$0$ cross-pairing
$\mathbf{16}_{-3}\!\otimes\!\overline{\mathbf{16}}_{+3}$, which is the
restriction of the Killing form. Hence the invariant bilinear is
one-dimensional (Schur, complex type), so unique up to scale. On the
compact form the Killing form is definite; restricted to $\mathfrak{m}$
it is positive-definite, giving the Euclidean target-space metric
$K|_{\mathfrak{m}}$ used in \autoref{sec:coset:main}.
\end{proof}

\begin{remark}[No Lorentzian content in the coset metric]\label{rmk:albert:nolorentz}
The coset metric $K|_{\mathfrak{m}}$ is positive-definite, signature
$(32,0)$: it carries \emph{zero} Lorentzian content. In particular there
is no $(4,28)$ Killing/coset signature, and the physical $(1,3)$ Minkowski
signature does \emph{not} fall out of the coset geometry. Lorentzian
signature enters only through the soldering postulate
\hyperref[post:soldering:P3]{P3}, which uses the indefinite cubic norm
$N$ restricted to $H_{2}(\mathbb{O})\cong\mathbb{R}^{1,9}$ on a
\emph{different} module (\autoref{sec:soldering:main}). The two forms
(definite $K|_{\mathfrak{m}}$, indefinite $N$) live on distinct spaces;
there is no tension.
\end{remark}

%%%%%%%%%%%%%%%%%%%%%%%%%%%%%%%%%%%%%%%%%%%%%%%%%%%%%%%%%%%%%%%%%%%%%%
\subsection{Summary of canonical structure}
\label{sec:albert:summary}
%%%%%%%%%%%%%%%%%%%%%%%%%%%%%%%%%%%%%%%%%%%%%%%%%%%%%%%%%%%%%%%%%%%%%%

For reference, \autoref{tab:albert:summary} collects the canonical
algebraic data established above; every downstream section uses these
values.

\begin{table}[ht]
\centering
\renewcommand{\arraystretch}{1.3}
\begin{tabular}{ll}
\hline
Object & Canonical value / property \\
\hline
$\dim_{\mathbb{R}} J_{3}(\mathbb{O})$ & $27$ \\
Trace form $Q=\operatorname{Tr}(X^{2})$ & positive-definite; $F_{4}$-invariant \emph{only} \\
Cubic norm $N=\det_{\mathbb{O}}X$ & indefinite; $E_{6}$-invariant up to character $\lambda(g)$ \\
Sharp map $X^{\#}=\nabla N$ & $E_{6}$-covariant; $X^{\#}=0\Leftrightarrow\operatorname{rank}\le1$ \\
$F_{4}=\operatorname{Aut}(J_{3}(\mathbb{O}))$ & compact, $\dim 52$; preserves $Q$ and $N$ \\
$E_{6(-26)}$ (orbit classifier) & non-compact, $\dim 78$; preserves $N$ only \\
compact $E_{6}$ / $E_{6(-14)}$ (dynamics) & $\dim 78$; max.\ compact $\operatorname{Spin}(10)\times\mathrm{U}(1)$ \\
Vacuum manifold $\mathcal{M}$ & $E_{6}/(\operatorname{Spin}(10)\times\mathrm{U}(1))=\mathrm{EIII}$, $\dim 32$ \\
Adjoint branching & $\mathbf{78}=\mathbf{45}_{0}\oplus\mathbf{1}_{0}\oplus\mathbf{16}_{-3}\oplus\overline{\mathbf{16}}_{+3}$ \\
Fundamental branching & $\mathbf{27}=\mathbf{1}_{+4}\oplus\mathbf{10}_{-2}\oplus\mathbf{16}_{+1}$ \\
Coset tangent $\mathfrak{m}$ & $\mathbf{16}_{-3}\oplus\overline{\mathbf{16}}_{+3}$ (32 real); no singlet \\
Coset metric $K|_{\mathfrak{m}}$ & positive-definite (Euclidean), unique up to scale \\
\hline
\end{tabular}
\caption{Canonical algebraic data of the Albert algebra and its
exceptional structure, as used throughout UCFT.}
\label{tab:albert:summary}
\end{table}

These facts are purely algebraic theorems, with the single
\emph{postulate} being the commitment to the compact real form for the
dynamics (\autoref{post:albert:realform}, part of
\hyperref[post:potential:P1]{P1}). The vacuum-selection theorem that
promotes \eqref{eq:albert:EIII} from ``a distinguished orbit'' to ``the
global minimum of $V$'' is proved in \autoref{sec:potential:main}; the
Lorentzian soldering that uses $N$ to install signature $(-,+,+,+)$ is in
\autoref{sec:soldering:main}; the gauge/clock-field dynamics on
$\mathfrak{m}$ are in \autoref{sec:coset:main}.
